\paragraph{\trSplitClosed{} A∞-\trCategory。}\label{sub:split-closed-categories}
若任何\trHomotopy\trIdempotent{}的抽象像都被 $\scr{A}$ 的某个\trObject{}%
\trQuasiRepresent，我们就称 $\scr{A}$ 是\trSplitClosed{}的。一个%
\trACategory{} $\scr{A}$ 的\trSplitClosure{}是一个对 $(\scr{B},\scr{F})$，由一个%
\trSplitClosed{} \trACategory{} $\scr{B}$ 连同一个\trCohomology%
\trFullyFaithful{}的\trFunctor{} $\scr{F}\colon\scr{A}\to\scr{B}$ 组成，使得在
$H^0(\scr{B})$ 中，每个\trObject{}都是 $H^0(\scr{A})$ 的某个\trObject{}的一个%
\trIdempotent{}\trEndomorphism{}的\trImage。

在陈述这些概念时，我们遵循了 §3 的一般模板；但事实上，它们的含义可以大大
简化。首先，$\scr{A}$ \emph{是\trSplitClosed{}的，当且仅当 $H^0(\scr{A})$
具有该性质}（按原先那个\trCohomology{}层面的意义）。具体说，假设 $H^0(\scr{A})$ 是%
\trSplitClosed{}的，并取一个\trHomotopy\trIdempotent{} $\wp$，使
$[\wp^1]=p$。由假设，$p$ 在 $H^0(\scr{A})$ 中有一个像 $Z$。利用%
\cref{thm:hom-abstract-image}容易证明：$Z$ 的\trYoneda{}模在 $H^0(\scr{Q})$ 中%
\trIsomorphic{}于 $\wp$ 的抽象像，从而给出一个\trQuasiRepresent{}元。反方向的
蕴涵类似，不过还要用到\cref{thm:homotopy-idempotent-lift}。一个直接的推论是：
$(\scr{B},\scr{F})$ 是 $\scr{A}$ 的一个\emph{\trSplitClosure}，当且仅当
$(H^0(\scr{B}),H^0(\scr{F}))$ 是 $H^0(\scr{A})$ 的一个\trSplitClosure。

\begin{lemma}\label{thm:split-closure-existence}
每个 \trACategory{}都有\trSplitClosure。进而，若 $(\scr{B},\scr{F})$ 与
$(\widetilde{\scr{B}},\widetilde{\scr{F}})$ 是 $\scr{A}$ 的两个这样的闭包，则存在%
\trQuasiequivalence{} $\scr{G}\colon\scr{B}\to\widetilde{\scr{B}}$，使得
$\scr{G}\circ\scr{F}$ 在 $H^0(\mathit{fun}(\scr{A},\scr{B}))$ 中与
$\widetilde{\scr{F}}$ \trIsomorphic。
\end{lemma}
\begin{proof}
把 $\Pi\scr{A}$ 记为 $\scr{Q}=\mathit{mod}(\scr{A})$ 的下述\trFull{}
\trAinfty-子\trCategory：它由全部在 $H^0(\scr{Q})$ 中与某个\trHomotopy%
\trIdempotent{} $\wp$ 的抽象像 $\scr{Z}$ \trIsomorphic{}的\trObject{}组成。我们
断言，$\Pi\scr{A}$ 包含与诸\trObject{} $Y$ 相关联的一切\trYoneda{}
\trAinfty-模 $\scr{Y}$。为看到这一点，取某个\trHomotopy\trIdempotent，使得
$[\wp^1]$ 是 $Y$ 的\trIdentityMorphism（由\cref{thm:homotopy-idempotent-lift}，
这是存在的）；据\cref{thm:hom-abstract-image}，$\wp$ 的抽象像将与 $Y$ 的%
\trYoneda{}模%
\trIsomorphic。有了这些，$(\Pi\scr{A},\ \scr{J}_{\scr{A}}\colon\scr{A}\to
\Pi\scr{A}\subset\scr{Q})$ 这一对便构成 $\scr{A}$ 的一个\trSplitClosure。诚然，正如上文所
观察到的，只需考察\trCohomology\trCategory{}即可，而在那里，
\cref{thm:homotopy-idempotent-lift}与\cref{thm:hom-abstract-image}的同一组合便
给出所要的结论。

对第二部分，设 $(\scr{B},\scr{F})$ 是 $\scr{A}$ 的任意一个\trSplitClosure。利用
$\scr{B}$ 的\trYonedaEmbedding{}与 \trAinfty-模的\trPullback，可以定义一个%
\trAFunctor
\begin{equation*}
\scr{G}\colon\scr{B}\xrightarrow{\ \scr{J}_{\scr{B}}\ }\mathit{mod}(\scr{B})
\xrightarrow{\ \scr{F}^*\ }\scr{Q}.
\end{equation*}
由于 $\scr{F}$ 是\trCohomology\trFullyFaithful{}的，$\scr{G}\circ\scr{F}$ 与
$\scr{J}_{\scr{A}}$ \trIsomorphic，见\eqref{eq:yoneda-naturality-diagram}。这
蕴含：把 $\scr{G}$ 限制到形如 $\scr{F}(Y)$ 的\trObject{}上时，它是%
\trCohomology\trFullyFaithful{}的；但按定义，$H^0(\scr{B})$ 的每个\trObject{}%
都是某个 $\scr{F}(Y)$ 的\trIdempotent{}\trEndomorphism{}的\trImage，因此
$\scr{G}$ 在整个 $\scr{B}$ 上是\trCohomology\trFullyFaithful{}的。类似的论证
表明，$\scr{G}$ 的\trImage{}含于 $\Pi\scr{A}$，并且 $\scr{G}$ 实际上是一个%
\trQuasiequivalence{} $\scr{B}\to\Pi\scr{A}$。当然，既然任意一个\trSplitClosure{}与
$\Pi\scr{A}$ 之间都有了\trQuasiequivalence，现在便可以借助%
\cref{thm:quasi-equivalence-homotopy-inverse}，在任意两个这样的闭包之间构造出%
\trQuasiequivalence{}来。
\end{proof}

\begin{lemma}\label{thm:split-closure-triangulated}
\trTriangulated{} \trACategory{}的\trSplitClosure{}仍是\trTriangulated{}的。
\end{lemma}
\begin{proof}
设 $\scr{A}$ 是一个\trTriangulated{} \trACategory。据%
\cref{thm:mod-categories-triangulated}，$\scr{Q}=\mathit{mod}(\scr{A})$ 是%
\trTriangulated{}的，而经由\trYonedaEmbedding，$H^0(\scr{A})$ 被等同于
$H^0(\scr{Q})$ 的一个\trTriangulated{}子\trCategory（经典意义下的）。此外，
$H^0(\Pi\scr{A})$ 是由那些与 $H^0(\scr{A})$ 中某个\trObject{}的一个%
\trDirectSummand{}\trIsomorphic{}的\trObject{}所组成的子\trCategory，因而据%
\cref{thm:summands-triangulated-subcategory}，它仍是 $H^0(\scr{Q})$ 的一个%
\trTriangulated{}子\trCategory。换言之，若在 $H(\scr{Q})$ 中有一个%
\trExactTriangle，而所涉及的\trObject{}中有两个落在 $\Pi\scr{A}$ 内，则第三个
亦然。据\cref{thm:preimage-exact-triangle}，这就蕴含 $\Pi\scr{A}$ 自身是一个%
\trTriangulated{} \trACategory。
\end{proof}

设 $\scr{A}$ 是任意一个 \trACategory；取它的一个\trTriangulated\trEnvelope，
然后再取后者的\trSplitClosure，并把所得结果记作 $\scr{B}$。例如，两次都取标准的选择
便得 $\scr{B}=\Pi(Tw\scr{A})$。我们把 $H^0(\scr{B})$ 称为\textbf{%
\trSplitClosed{}\trDerivedCategory}（split-closed derived category），并
记作 $D^\pi(\scr{A})$。

与\cref{sub:generators}中的讨论平行，我们还应提一个相关的概念。设 $\scr{B}$
是\trSplitClosed{} \trTriangulated{} \trACategory，而 $\scr{A}\subset\scr{B}$
是一个不空的\trFull{} \trAinfty-子\trCategory。考虑最小的\trFull{}子\trCategory{}%
$\widetilde{\scr{B}}\subset\scr{B}$：它包含 $\scr{A}$，对 $H^0(\scr{B})$ 中的%
\trIsomorphism{}封闭，而且其自身是\trTriangulated{}且\trSplitClosed{}的。
我们称之为 $\scr{B}$ 中由 $\scr{A}$ \trSplitGeneration{}的子\trCategory。具体
说来，从 $\Ob\scr{A}$ 出发，作出\trMappingCone{}与\trSuspension，把这一过程
重复任意多次，最后再取先前构造出的诸\trObject{}在一切\trIdempotent{}%
\trEndomorphism{}下的\trImage，便得到 $\Ob\widetilde{\scr{B}}$。若
$\widetilde{\scr{B}}=\scr{B}$，我们就说 $\scr{A}$ \emph{\trSplitGeneration}
$\scr{B}$（或者说，它的\trObject{}是 $\scr{B}$ 的\trSplitGeneration{}元）。
例如，$\Pi(Tw\scr{A})$ 由 $\scr{A}$ 的\trObject{}所\trSplitGeneration。作为
\cref{thm:triangulated-envelope-existence}与\cref{thm:split-closure-existence}
中唯一性结果的一个直接推论，我们得到：

\begin{corollary}\label{thm:split-generated-equivalence}
设 $\scr{B}$ 是\trSplitClosed{} \trTriangulated{} \trACategory，并设
$\scr{A}\subset\scr{B}$ 是一个\trSplitGeneration{}它的\trFull{}子\trCategory。
那么存在\trQuasiequivalence{} $\Pi(Tw\scr{A})\to\scr{B}$，它诱导出诸%
\trTriangulated{}\trCategory{}之间的一个\trEquivalence{}%
$H^0(\scr{B})\cong D^\pi(\scr{A})$。
\noqed
\end{corollary}

\begin{remark}\label{thm:unbounded-twisted-alternative}
\trSplitClosure $\Pi\scr{A}$ 有一点在美观上不讨人喜欢：它并不把 $\scr{A}$ 自身作为%
\trFull{}子\trCategory{}包含在内（当然，它确实包含一个与 $\scr{A}$ %
\trQuasiisomorphic{}的子\trCategory）。下面我们概述另一种构造，它是无界%
\trTwistedComplex{}的一个变体。这一构造更为具体，也更贴近最初的灵感，但至少
就这里给出的版本而言，它要求\trStrictlyUnital{}性。
\end{remark}

设 $\scr{A}$ 是一个\trStrictlyUnital{}的 \trACategory。我们将假定它有一个零%
\trObject{} $O$，使得 $\hom_{\scr{A}}(O,X)$ 与 $\hom_{\scr{A}}(X,O)$ 都是平凡
群。按如下方式定义另一个 \trACategory{} $\Sigma^-\scr{A}$。$\Sigma^-\scr{A}$ 的%
\trObject{}是形式的无穷序列 $X=(X^0,X^{-1},X^{-2},\ldots)$，其中
$X^i\in\Ob\scr{A}$。\trMorphism{} $a\in\hom^k_{\Sigma^-\scr{A}}(X_0,X_1)$ 是一个
无穷矩阵
\begin{equation*}
a=(a^{ji}),\qquad a^{ji}\in\hom^{k+j-i}_{\scr{A}}\bigl(X_0^i,X_1^j\bigr),
\end{equation*}
它是“本质下三角”的，意即：存在某个 $N$（依赖于 $a$），使得当 $j<i-N$ 时
$a^{ji}=0$。\trComposition{}映射用矩阵乘法来定义，方式与\trAdditive{}扩大处
相同。注意，每个 $X^i$ 分块都被看作\trSuspension{}了 $i$ 次，这就带来一些正负号。事实上，
为使下面的公式简化，我们采用\cref{thm:shifted-multiplicity-spaces}中修正过的
正负号约定。

考虑对 $(X,\delta_X)$：其中 $X\in\Ob\Sigma^-\scr{A}$，而
$\delta_X\in\hom^1_{\Sigma^-\scr{A}}(X,X)$ 在 $\delta_X^{ji}=0$（对 $j\le i$）的
意义下是严格下三角的，并且满足\eqref{eq:generalized-maurer-cartan}。我们定义
一个 \trACategory{} $Tw^-\scr{A}$：其\trObject{}就是这种对，而其\trMorphism{}%
则如\eqref{eq:twisted-composition} 那样被诸\trDifferential{}所\trDeformation。
可以把 $\scr{A}$ 等同于 $Tw^-\scr{A}$ 的由序列
$(X^0,X^{-1}=O,X^{-2}=O,\ldots)$ 组成的\trFull{} \trAinfty-子\trCategory。

给定了 $\Ob\scr{A}$ 中\trObject{} $Y$ 的一个\trHomotopy\trIdempotent{} $\wp$，
我们如下定义 $Tw^-\scr{A}$ 的一个\trObject{} $(Z,\delta_Z)$：
\begin{equation}\label{eq:homotopy-idempotent-object}
\begin{aligned}
&Z^i=Y\quad\text{对一切 }i,\\
&\delta_Z^{ji}=
\begin{cases}
\wp^1 & \text{若 } i \text{ 为偶数，且 } j=i+1\le 0,\\
\wp^1-e_Y & \text{若 } i \text{ 为奇数，且 } j=i+1\le 0,\\
\wp^{j-i} & \text{若 } i+1<j\le 0,\\
0 & \text{在其余一切情形。}
\end{cases}
\end{aligned}
\end{equation}
广义 \trMaurerCartan{}方程\trEquivalence{}于\eqref{eq:homotopy-idempotent-equations}。定义
\begin{equation*}
\begin{aligned}
&r\in\hom^0_{Tw^-\scr{A}}(Y,Z),\quad &&r^{00}=e_Y,\\
&k\in\hom^0_{Tw^-\scr{A}}(Z,Y),\quad &&k^{0i}=\wp^{1-i},\\
&h\in\hom^{-1}_{Tw^-\scr{A}}(Z,Z),\quad &&h^{i-1,i}=(-1)^i e_Y,
\end{aligned}
\end{equation*}
其中凡未提及的\trCoefficients{}均为 $0$。于是
\begin{equation*}
\begin{aligned}
&\mu^1_{Tw^-\scr{A}}(r)=0, &&\mu^1_{Tw^-\scr{A}}(k)=0,\\
&\mu^2_{Tw^-\scr{A}}(k,r)=\wp^1,
&&\mu^2_{Tw^-\scr{A}}(k,r)-\mu^1_{Tw^-\scr{A}}(h)=e_Z.
\end{aligned}
\end{equation*}

这就意味着，在\trCohomology\trCategory{}中，$Z$ 是 $Y$ 的\trIdempotent{}%
\trEndomorphism{} $[\wp^1]$ 的\trImage。于是，若取 $Tw^-\scr{A}$ 的这样一个%
\trFull{} \trAinfty-子\trCategory：它由 $\scr{A}$ 的全部\trObject{}连同
\eqref{eq:homotopy-idempotent-object} 型的那些\trObject{}组成，则这是一个%
\trSplitClosure。为了与前面的做法建立联系，请注意与 $Z$ 相关联的%
\trYoneda{} $\scr{A}$-模恰好就是抽象像 $\scr{Z}$。
